\appendix{Представление графического материала}

Графический материал, выполненный на отдельных листах,
изображен на рисунках А.1--А.\arabic{числоПлакатов}.
\setcounter{числоПлакатов}{0}

\renewcommand{\thefigure}{А.\arabic{figure}} % шаблон номера для плакатов

\begin{landscape}

\begin{плакат}
    
\includegraphics[width=0.82\linewidth]{Ramka_VKR (1).png}
    \заголовок{Сведения о ВКРБ}
    \label{pl1:image}      
\end{плакат}

\begin{плакат}
    \includegraphics[width=0.82\linewidth]{tasks (2).png}
    \заголовок{Цель и задачи разработки}
    \label{pl2:image}      
\end{плакат}

\begin{плакат}
    \includegraphics[width=0.82\linewidth]{concurent (3).png}
    \заголовок{Сравнение существующих онлайн-площадок для аренды вещей}
    \label{pl3:image}      
\end{плакат}

\begin{плакат}
    \includegraphics[width=0.82\linewidth]{components (4).png}
    \заголовок{Компоненты онлайн-площадки для аренды вещей}
    \label{pl4:image}      
\end{плакат}

\begin{плакат}
    \includegraphics[width=0.82\linewidth]{diaprec (5).png}
    \заголовок{Диаграмма прецедентов}
    \label{pl5:image}      
\end{плакат}

\begin{плакат}
    \includegraphics[width=0.82\linewidth]{diacomp (6).png}
    \заголовок{Диаграмма компонентов}
    \label{pl6:image}      
\end{плакат}

\begin{плакат}
    \includegraphics[width=0.82\linewidth]{api-table (7).png}
    \заголовок{Основные HTTP-маршруты API}
    \label{pl7:image}      
\end{плакат}

\begin{плакат}
    \includegraphics[width=0.82\linewidth]{test (8).png}
    \заголовок{Тестирование модулей. Часть 1}
    \label{pl8:image}      
\end{плакат}

\begin{плакат}
    \includegraphics[width=0.82\linewidth]{test (9).png}
    \заголовок{Тестирование модулей. Часть 2}
    \label{pl9:image}      
\end{плакат}

\begin{плакат}
    \includegraphics[width=0.82\linewidth]{end (10).png}
    \заголовок{Заключение}
    \label{pl10:image}      
\end{плакат}

\end{landscape}
